@@M0@@) @@M1@@ are linearly dependent over @@M2@@.

@@M3@@) @@M4@@ are linearly independent over @@M5@@.

\sgaStableTarget{sga6.e06.sec.1.2}1.2.
For convenience, we shall recall here a few classical definitions and properties of graded Modules.

Let @@M6@@ be a scheme, @@M7@@ a graded @@M8@@-Algebra, @@M9@@ a graded @@M10@@-Module, For every graded @@M11@@-Module @@M12@@, we denote as usual by @@M13@@ the graded @@M14@@-Module defined by:
@@M15@@
We say that @@M16@@ is a locally free graded @@M17@@-Module (resp. locally free of finite type) if every point of @@M18@@ has an open neighborhood @@M19@@ such that @@M20@@ is isomorphic, as a graded @@M21@@-Module, to a direct sum of Modules of the form @@M22@@ (resp. a finite direct sum).

We say that @@M23@@ is a graded @@M24@@-Module of finite presentation if every point of @@M25@@ has an open neighborhood @@M26@@ such that there is an exact sequence @@M27@@, where @@M28@@ and @@M29@@ are free graded @@M30@@-Modules of finite type, and the homomorphisms are compatible with the gradings.

1.2.1. For every integer @@M31@@, we can define a canonical homomorphism of graded @@M32@@-Modules:
\sgaStableTarget{sga6.e06.eq.1-2-1}@@M33@@
by sending @@M34@@, where @@M35@@ is a section of @@M36@@, and @@M37@@ a section of @@M38@@, to the element @@M39@@ of @@M40@@.

1.2.2. We say that a graded @@M41@@-Module @@M42@@ is free of finite type and homogeneous of degree @@M43@@ if it is isomorphic as a graded @@M44@@-Module to @@M45@@ for some integer @@M46@@.

